\section{Conclusion}
\label{sec:conclusion}

This report presented a complete, from-scratch implementation of a
decoder-only Transformer (GPT-style) language model in PyTorch. Every
core component---token embeddings, positional embeddings, causal
multi-head self-attention, Pre-LN Transformer blocks, GELU feed-forward
networks, and autoregressive generation---was built without relying on
high-level wrappers such as \texttt{nn.TransformerDecoder}. This
constraint turned out to be the most valuable part of the exercise: it
forced us to reason about tensor shapes, masking semantics, and gradient
flow at a level that high-level abstractions normally hide.

\subsection{Summary of Findings}
Training ran for 5000 steps and completed in \textbf{3.83 minutes} on a
laptop GPU, converging to a final training loss of $1.613$ and a
validation loss of $1.545$. In perplexity terms, the model achieves
\textbf{4.69} on the held-out split---a meaningful improvement over the
character-level bigram baseline of roughly $3.4$ on this corpus, though
still short of the $\approx 1.8$--$2.0$ that a fully-trained character
LSTM would reach. Given that the model has only $0.82$~M parameters and
sees roughly $1.2$ epochs over the corpus, this operating point is
exactly what the theory predicts: the model captures local character
statistics and short-range syntax, but lacks the training exposure
required to fit longer-range dependencies.

\noindent The three optimisation choices highlighted in the assignment all behaved
as expected. Mixed precision (bf16 via \texttt{autocast}) preserved
numerical stability throughout the run, with no NaNs, no loss spikes,
and no evidence of gradient underflow. The linear-warmup-cosine-decay
schedule produced the smoothest part of the training curve in its final
phase, with validation loss descending monotonically from $1.559$ at
step $4250$ to $1.545$ at step $5000$ as the learning rate decayed.
Selective weight decay---applied to $18$ tensors with
$\dim(\theta) \geq 2$ but excluded from $34$ one-dimensional
tensors---matched the standard GPT recipe and required no further
tuning.

\noindent The inference loop supports greedy decoding, temperature scaling, and
top-$k$ sampling. The comparison across five strategies from the same
prompt was the most visually satisfying result of the project: greedy
collapsed into a repetitive loop within fifteen tokens, temperature
$0.8$ with top-$k = 40$ produced fluent, varied Shakespeare-shaped
text, and temperature $1.5$ degenerated into gibberish. The
coherence-diversity trade-off that the literature describes
\parencite{holtzman2020curious} showed up exactly as advertised, and
being able to observe it directly---rather than take it on faith---was
the point of building the model from scratch in the first place.

\subsection{What This Project Taught Me}
Beyond the mechanics, three things stand out.

\noindent First, the causal mask is not a minor implementation detail; it is the
mechanism that makes autoregressive training mathematically equivalent
to parallel decoding. Until you have written the
\texttt{masked\_fill} call yourself and watched the lower-triangular
attention heatmaps emerge, it is easy to treat the mask as boilerplate.
It is not.

\noindent Second, mixed precision is not just a speed optimisation. Working
through the interaction between \texttt{autocast}, \texttt{GradScaler},
gradient clipping, and unscaling is what clarified for me why bf16 and
fp16 have different tooling requirements---and why getting the order of
operations wrong silently destroys the clip threshold.

\noindent Third, and perhaps most importantly, the small implementation choices
that textbooks often skim over---Pre-LN vs.\ Post-LN, whether to apply
weight decay to LayerNorm gains, whether to tie the embedding and LM
head weights---turn out to matter more than I expected for training
stability and final model quality.

\subsection{Future Work}
Several extensions are natural next steps:

\begin{itemize}
    \item \textbf{KV-cache.} The current generation loop recomputes the
          entire context at every step, giving $O(T^2)$ cost per
          generated token. A KV-cache would reduce this to $O(T)$ and
          make long-form generation practical.
    \item \textbf{Top-$p$ (nucleus) sampling.} Adapting $k$ to the shape
          of each token distribution would likely outperform fixed
          top-$k$, especially for prompts whose next-token entropy
          varies sharply.
    \item \textbf{Repetition penalties.} A small penalty on already-seen
          tokens would directly address the looping behaviour observed
          under greedy decoding.
    \item \textbf{Rotary Positional Embeddings (RoPE).} Replacing
          learned/sinusoidal embeddings with RoPE would scale better to
          longer contexts and is now standard in most modern
          open-weight LLMs \parencite{su2021roformer}.
    \item \textbf{Multi-GPU training.} Distributed Data Parallel or
          Fully Sharded Data Parallel would allow the model to grow
          past the single-GPU memory budget.
    \item \textbf{Parameter-efficient fine-tuning.} LoRA or a similar
          adapter-based method would make it possible to adapt the
          pretrained checkpoint to a different corpus without full
          retraining \parencite{hu2021lora}.
\end{itemize}

\noindent Looking back, the value of this assignment was not in producing
state-of-the-art text---Tiny Shakespeare at $0.82$~M parameters was
never going to compete with a modern LLM. The value was in removing
every layer of abstraction between the equations on the page and the
tensors in memory. Having done that, I now read papers on Transformers
differently: I can see where the implementation complexity actually
lives, and I understand which architectural choices are load-bearing
versus which are cosmetic. For a $4$-minute training run, that is a
remarkably good return on investment.